% Minimalist light Beamer styling for academic and engineering slides.
% Loaded by _quarto.yml; individual notes do not need raw LaTeX.

\usepackage{sourceserifpro}
\usepackage{microtype}
\usepackage{etoolbox}
\usepackage{graphicx}
\usepackage{fvextra}
\usepackage{tikz}
\usepackage{xcolor}

% pdfLaTeX compatibility for invisible Obsidian characters and common
% Unicode mathematics pasted into Markdown equations.
\DeclareUnicodeCharacter{200B}{}
\DeclareUnicodeCharacter{00A0}{~}
\DeclareUnicodeCharacter{03B8}{\ensuremath{\theta}}
\DeclareUnicodeCharacter{03BB}{\ensuremath{\lambda}}
\DeclareUnicodeCharacter{03C0}{\ensuremath{\pi}}
\DeclareUnicodeCharacter{03C3}{\ensuremath{\sigma}}
\DeclareUnicodeCharacter{03C4}{\ensuremath{\tau}}
\DeclareUnicodeCharacter{2113}{\ensuremath{\ell}}
\DeclareUnicodeCharacter{2192}{\ensuremath{\rightarrow}}
\DeclareUnicodeCharacter{2212}{\ensuremath{-}}
\DeclareUnicodeCharacter{2223}{\ensuremath{\mid}}
\DeclareUnicodeCharacter{2264}{\ensuremath{\leq}}
\DeclareUnicodeCharacter{2265}{\ensuremath{\geq}}
\DeclareUnicodeCharacter{2026}{\ensuremath{\ldots}}

% Source Serif text with Computer Modern mathematics.
% Source Serif changes the text family only; the explicit symbol
% declarations below retain Computer Modern for equations.
\SetSymbolFont{operators}{normal}{OT1}{cmr}{m}{n}
\SetSymbolFont{letters}{normal}{OML}{cmm}{m}{it}
\SetSymbolFont{symbols}{normal}{OMS}{cmsy}{m}{n}
\SetSymbolFont{largesymbols}{normal}{OMX}{cmex}{m}{n}
\SetSymbolFont{operators}{bold}{OT1}{cmr}{bx}{n}
\SetSymbolFont{letters}{bold}{OML}{cmm}{b}{it}
\SetSymbolFont{symbols}{bold}{OMS}{cmsy}{b}{n}
\usefonttheme{serif}

\definecolor{MinimalInk}{HTML}{202124}
\definecolor{MinimalMuted}{HTML}{70757A}
\definecolor{MinimalFaint}{HTML}{A7AAAD}
\definecolor{MinimalLink}{HTML}{426B86}
\definecolor{MinimalCode}{HTML}{F5F6F7}
\definecolor{MinimalRule}{HTML}{E5E7E9}

\setbeamercolor{background canvas}{bg=white}
\setbeamercolor{normal text}{fg=MinimalInk,bg=white}
\setbeamercolor{structure}{fg=MinimalInk}
\setbeamercolor{title}{fg=MinimalInk}
\setbeamercolor{subtitle}{fg=MinimalInk}
\setbeamercolor{author}{fg=MinimalInk}
\setbeamercolor{institute}{fg=MinimalInk}
\setbeamercolor{date}{fg=MinimalInk}
\setbeamercolor{frametitle}{fg=MinimalInk,bg=white}
\setbeamercolor{page number in head/foot}{fg=MinimalFaint,bg=white}
\setbeamercolor{block title}{fg=MinimalInk,bg=white}
\setbeamercolor{block body}{fg=MinimalInk,bg=white}
\setbeamercolor{alerted text}{fg=MinimalInk}
\setbeamercolor{example text}{fg=MinimalInk}

\setbeamerfont{normal text}{family=\rmfamily,size=\fontsize{8.6}{12}}
\setbeamerfont{itemize/enumerate body}{parent=normal text}
\setbeamerfont{itemize/enumerate subbody}{parent=normal text}
\setbeamerfont{title}{family=\rmfamily,series=\mdseries,size=\fontsize{15}{18}}
\setbeamerfont{subtitle}{family=\rmfamily,series=\mdseries,size=\fontsize{11}{13.5}}
\setbeamerfont{author}{family=\rmfamily,size=\fontsize{9.5}{12}}
\setbeamerfont{institute}{family=\rmfamily,size=\fontsize{8.5}{11}}
\setbeamerfont{date}{family=\rmfamily,size=\fontsize{8.5}{10.5}}
\setbeamerfont{frametitle}{family=\rmfamily,series=\mdseries,size=\fontsize{13.5}{16.5}}
\setbeamerfont{block title}{family=\rmfamily,series=\bfseries,size=\normalsize}
\setbeamerfont{caption}{family=\rmfamily,series=\mdseries,size=\fontsize{6.5}{8}}
\setbeamerfont{page number in head/foot}{family=\rmfamily,size=\scriptsize}

\setbeamersize{text margin left=0.72cm,text margin right=0.72cm}
\setlength{\parindent}{0pt}
\setlength{\parskip}{0.42em}
\setlength{\columnsep}{0.65cm}
% Align first-level bullet markers with the body margin and keep nesting tight.
\setlength{\leftmargini}{0.90em}
\setlength{\leftmarginii}{1.05em}

\setbeamertemplate{navigation symbols}{}
\setbeamertemplate{headline}{}
\setbeamertemplate{frametitle continuation}{}
\setbeamertemplate{section in toc}{}
\setbeamertemplate{subsection in toc}{}
\setbeamertemplate{blocks}[default]
\setbeamertemplate{caption}{\raggedright\insertcaption\par}
\setbeamertemplate{caption label separator}{}
\renewcommand{\figurename}{}
\renewcommand{\tablename}{}
\setbeamertemplate{itemize item}{\raisebox{0.16ex}{\scriptsize$\bullet$}}
\setbeamertemplate{itemize subitem}{\raisebox{0.12ex}{\tiny$\circ$}}
\setbeamertemplate{enumerate item}{\insertenumlabel.}

\setbeamertemplate{frametitle}{%
  \nointerlineskip
  % Draw the title at the top without consuming vertical frame space. The
  % body is therefore centered against the whole slide, not the region left
  % over below the title.
  \vbox to 0pt{%
    \vspace*{0.464cm}%
    \begin{beamercolorbox}[
      wd=\paperwidth,
      sep=0.28cm,
      leftskip=0.44cm,
      rightskip=0.44cm
    ]{frametitle}%
      \usebeamerfont{frametitle}\insertframetitle\strut
    \end{beamercolorbox}%
    \vss
  }%
}

\setbeamertemplate{footline}{%
  \leavevmode
  \hbox{%
    \begin{beamercolorbox}[wd=\paperwidth,ht=2.5ex,dp=1.15ex,right]{page number in head/foot}%
      \ifnum\value{framenumber}>0\relax
        \usebeamerfont{page number in head/foot}\insertframenumber\hspace*{0.65cm}%
      \fi
    \end{beamercolorbox}%
  }%
}

\makeatletter
\newcommand{\obsbeambodyspacing}{%
  \fontsize{8.6}{12}\selectfont%
  \setlength{\parskip}{0.38em}%
  \renewcommand{\arraystretch}{1.07}%
  \setlength{\abovedisplayskip}{0.57em}%
  \setlength{\belowdisplayskip}{0.57em}%
  \setlength{\abovedisplayshortskip}{0.28em}%
  \setlength{\belowdisplayshortskip}{0.38em}%
}
\newcommand{\obsbeamtightlist}{%
  \ifnum\@itemdepth>1\relax
    \setlength{\itemsep}{0.09em}%
    \setlength{\parskip}{0pt}%
    \setlength{\parsep}{0pt}%
  \else
    \setlength{\itemsep}{0.17em}%
    \setlength{\parskip}{0pt}%
    \setlength{\parsep}{0pt}%
  \fi
}
% Equal stretch plus a small fixed-bound difference places the visible body
% 0.25cm below the mathematical page center for a gentler optical balance.
\define@key{beamerframe}{s}[true]{%
  \beamer@frametopskip=1.50cm plus 1fill\relax%
  \beamer@framebottomskip=1.00cm plus 1fill\relax%
  \beamer@frametopskipautobreak=\beamer@frametopskip\relax%
  \beamer@framebottomskipautobreak=\beamer@framebottomskip\relax%
  \obsbeambodyspacing%
}
% Media frames use the same fixed safe area, but figures and column layouts
% are indivisible in Beamer and must not trigger an empty continuation page.
\define@key{beamerframe}{m}[true]{%
  \beamer@frametopskip=1.40cm plus 1fill\relax%
  \beamer@framebottomskip=1.10cm plus 1fill\relax%
  \beamer@frametopskipautobreak=\beamer@frametopskip\relax%
  \beamer@framebottomskipautobreak=\beamer@framebottomskip\relax%
  \obsbeambodyspacing%
}
\newlength{\obsbeam@titleoffset}
\newlength{\obsbeam@titlegap}
\newcommand{\obsbeam@authorblock}{%
  \ifx\insertauthor\@empty\else
    {\usebeamerfont{author}\usebeamercolor[fg]{author}\insertauthor\par}%
  \fi
  \ifx\insertinstitute\@empty\else
    \vspace{0.12cm}%
    {\usebeamerfont{institute}\usebeamercolor[fg]{institute}\insertinstitute\par}%
  \fi
  \ifx\insertdate\@empty\else
    \vspace{0.18cm}%
    {\usebeamerfont{date}\usebeamercolor[fg]{date}\insertdate\par}%
  \fi
}
\newcommand{\obsbeam@authorstack}{%
  \begin{tabular}{@{}c@{}}
    \ifx\insertauthor\@empty\else
      {\usebeamerfont{author}\usebeamercolor[fg]{author}\insertauthor}%
    \fi
    \ifx\insertinstitute\@empty\else
      \\[0.12cm]
      {\usebeamerfont{institute}\usebeamercolor[fg]{institute}\insertinstitute}%
    \fi
    \ifx\insertdate\@empty\else
      \\[0.18cm]
      {\usebeamerfont{date}\usebeamercolor[fg]{date}\insertdate}%
    \fi
  \end{tabular}%
}
\setbeamertemplate{title page}{%
  \setlength{\obsbeam@titleoffset}{1.60cm}%
  \setlength{\obsbeam@titlegap}{1.15cm}%
  \begin{minipage}[c][0.76\paperheight][t]{\textwidth}
    \centering
    \vspace*{\obsbeam@titleoffset}
    \begin{minipage}{0.72\textwidth}
      \centering
      {\usebeamerfont{title}\usebeamercolor[fg]{title}\inserttitle\par}
    \end{minipage}\par
    \ifx\insertsubtitle\@empty\else
      \vspace{\obsbeam@titlegap}
      {\usebeamerfont{subtitle}\usebeamercolor[fg]{subtitle}\insertsubtitle\par}
    \fi
    \vspace{\obsbeam@titlegap}
    \ifx\inserttitlegraphic\@empty
      \obsbeam@authorblock
    \else
      \makebox[\textwidth][c]{%
        \obsbeam@authorstack
        \hspace{0.253cm}%
        \raisebox{\dimexpr-.5\height+0.22cm\relax}{\inserttitlegraphic}%
      }%
    \fi
  \end{minipage}
  \addtocounter{framenumber}{-1}%
}
\makeatother

\addtobeamertemplate{itemize/enumerate body begin}{%
  \setlength{\baselineskip}{14.5pt}% wrapped lines within one item
  \setlength{\itemsep}{0.17em}%
  \setlength{\topsep}{0.38em}%
  \setlength{\parsep}{0.38em}%
}{}
\addtobeamertemplate{itemize/enumerate subbody begin}{%
  \setlength{\baselineskip}{14.5pt}% wrapped lines within one nested item
  \setlength{\itemsep}{0.09em}%
  \setlength{\topsep}{0.24em}%
  \setlength{\parsep}{0.24em}%
}{}
\AtBeginEnvironment{figure}{\renewcommand{\figurename}{}}
\AtBeginEnvironment{table}{\renewcommand{\tablename}{}}

% Predictable image behavior: fill available width, never distort, and never
% exceed a comfortable content height unless the Markdown requests it.
\setkeys{Gin}{width=\linewidth,totalheight=0.72\textheight,keepaspectratio}

% Embedded audio/video for Firefox's PDF.js media annotation player. The
% media bytes live inside the PDF; the visible box remains a quiet poster in
% viewers that do not support RichMedia annotations.
\makeatletter
\newcount\obsbeam@mediacount
\newdimen\obsbeam@mediawidth
\newdimen\obsbeam@mediaheight
\newdimen\obsbeam@mediamaxheight
\newbox\obsbeam@mediaposter
\newcommand{\obsbeammedia}[8]{%
  \begingroup
  \global\advance\obsbeam@mediacount by 1\relax
  \obsbeam@mediawidth=#4\relax
  \obsbeam@mediamaxheight=#5\relax
  \obsbeam@mediaheight=\dimexpr\obsbeam@mediawidth*#6/#7\relax
  \ifdim\obsbeam@mediaheight>\obsbeam@mediamaxheight
    \obsbeam@mediaheight=\obsbeam@mediamaxheight
    \obsbeam@mediawidth=\dimexpr\obsbeam@mediaheight*#7/#6\relax
  \fi
  \immediate\pdfextension obj stream
    attr {/Type /EmbeddedFile /Subtype /\pdf@escapename{#8}}
    file {#1}%
  \edef\obsbeam@streamobj{%
    \the\numexpr\pdffeedback lastobj\relax\space 0 R}%
  \immediate\pdfextension obj {<<
    /Type /Filespec
    /F <#2>
    /UF <#2>
    /EF << /F \obsbeam@streamobj /UF \obsbeam@streamobj >>
  >>}%
  \edef\obsbeam@filespecobj{%
    \the\numexpr\pdffeedback lastobj\relax\space 0 R}%
  \immediate\pdfextension obj {<<
    /Type /RichMediaInstance
    /Subtype /#3
    /Asset \obsbeam@filespecobj
  >>}%
  \edef\obsbeam@instanceobj{%
    \the\numexpr\pdffeedback lastobj\relax\space 0 R}%
  \immediate\pdfextension obj {<<
    /Type /RichMediaConfiguration
    /Subtype /#3
    /Instances [\obsbeam@instanceobj]
  >>}%
  \edef\obsbeam@configobj{%
    \the\numexpr\pdffeedback lastobj\relax\space 0 R}%
  \immediate\pdfextension obj {<<
    /Type /RichMediaContent
    /Assets << /Names [<#2> \obsbeam@filespecobj] >>
    /Configurations [\obsbeam@configobj]
  >>}%
  \edef\obsbeam@contentobj{%
    \the\numexpr\pdffeedback lastobj\relax\space 0 R}%
  \setbox\obsbeam@mediaposter=\hbox{%
    \begin{tikzpicture}[baseline=(current bounding box.south)]
      \path[use as bounding box]
        (0,0) rectangle (\obsbeam@mediawidth,\obsbeam@mediaheight);
      \fill[MinimalCode]
        (0,0) rectangle (\obsbeam@mediawidth,\obsbeam@mediaheight);
      \draw[MinimalRule,line width=0.45pt]
        (0,0) rectangle (\obsbeam@mediawidth,\obsbeam@mediaheight);
      \node[
        circle,
        fill=white,
        draw=MinimalRule,
        minimum size=1.05cm,
        inner sep=0pt
      ] at (.5\obsbeam@mediawidth,.5\obsbeam@mediaheight)
        {\color{MinimalInk}\large$\blacktriangleright$};
    \end{tikzpicture}%
  }%
  \leavevmode\hbox{%
    \rlap{%
      \pdfextension annot
        width \obsbeam@mediawidth
        height \obsbeam@mediaheight
        depth 0pt {%
          /Subtype /RichMedia
          /F 4
          /Border [0 0 0]
          /BS << /W 0 >>
          /NM (obsbeam-media-\the\obsbeam@mediacount)
          /Contents <#2>
          /RichMediaContent \obsbeam@contentobj
        }%
    }%
    \box\obsbeam@mediaposter
  }%
  \endgroup
}
\makeatother

% Quiet code treatment. Pandoc supplies syntax colors; this controls fit.
\AtBeginDocument{%
  \providecommand{\tightlist}{}%
  \renewcommand{\tightlist}{\obsbeamtightlist}%
  \fvset{breaklines=true,breakanywhere=true,fontsize=\small,baselinestretch=1.5,commandchars=\\\{\}}
  % Quarto/Pandoc emits real figure captions. Beamer's default template adds
  % "Figure" before them; enforce caption text only after all theme setup.
  \setbeamertemplate{caption}{\raggedright\insertcaption\par}%
  \setbeamertemplate{caption label separator}{}%
  \renewcommand{\figurename}{}%
  \renewcommand{\tablename}{}%
}
\AtBeginEnvironment{Shaded}{\small}

\hypersetup{
  colorlinks=true,
  linkcolor=MinimalInk,
  urlcolor=MinimalLink,
  citecolor=MinimalInk
}

% Mark external Markdown links automatically. The arrow is part of the link,
% so notes need no special syntax and the whole label remains clickable.
\AtBeginDocument{%
  \let\MinimalOriginalHref\href
  \renewcommand{\href}[2]{%
    \MinimalOriginalHref{#1}{%
      \textcolor{MinimalLink}{%
        #2\,\raisebox{0.12ex}{\scriptsize$\nearrow$}%
      }%
    }%
  }%
}
